% Options for packages loaded elsewhere
\PassOptionsToPackage{unicode}{hyperref}
\PassOptionsToPackage{hyphens}{url}
%
\documentclass[
  11pt,
]{article}
\usepackage{amsmath,amssymb}
\usepackage{iftex}
\ifPDFTeX
  \usepackage[T1]{fontenc}
  \usepackage[utf8]{inputenc}
  \usepackage{textcomp} % provide euro and other symbols
\else % if luatex or xetex
  \usepackage{unicode-math} % this also loads fontspec
  \defaultfontfeatures{Scale=MatchLowercase}
  \defaultfontfeatures[\rmfamily]{Ligatures=TeX,Scale=1}
\fi
\usepackage{lmodern}
\ifPDFTeX\else
  % xetex/luatex font selection
\fi
% Use upquote if available, for straight quotes in verbatim environments
\IfFileExists{upquote.sty}{\usepackage{upquote}}{}
\IfFileExists{microtype.sty}{% use microtype if available
  \usepackage[]{microtype}
  \UseMicrotypeSet[protrusion]{basicmath} % disable protrusion for tt fonts
}{}
\makeatletter
\@ifundefined{KOMAClassName}{% if non-KOMA class
  \IfFileExists{parskip.sty}{%
    \usepackage{parskip}
  }{% else
    \setlength{\parindent}{0pt}
    \setlength{\parskip}{6pt plus 2pt minus 1pt}}
}{% if KOMA class
  \KOMAoptions{parskip=half}}
\makeatother
\usepackage{xcolor}
\usepackage[margin=2.5cm]{geometry}
\usepackage{longtable,booktabs,array}
\usepackage{calc} % for calculating minipage widths
% Correct order of tables after \paragraph or \subparagraph
\usepackage{etoolbox}
\makeatletter
\patchcmd\longtable{\par}{\if@noskipsec\mbox{}\fi\par}{}{}
\makeatother
% Allow footnotes in longtable head/foot
\IfFileExists{footnotehyper.sty}{\usepackage{footnotehyper}}{\usepackage{footnote}}
\makesavenoteenv{longtable}
\setlength{\emergencystretch}{3em} % prevent overfull lines
\providecommand{\tightlist}{%
  \setlength{\itemsep}{0pt}\setlength{\parskip}{0pt}}
\setcounter{secnumdepth}{-\maxdimen} % remove section numbering
\ifLuaTeX
  \usepackage{selnolig}  % disable illegal ligatures
\fi
\IfFileExists{bookmark.sty}{\usepackage{bookmark}}{\usepackage{hyperref}}
\IfFileExists{xurl.sty}{\usepackage{xurl}}{} % add URL line breaks if available
\urlstyle{same}
\hypersetup{
  pdftitle={Bipartite graphs of diameter 7 from the split Cayley hexagon: level graphs, surgery and hubs},
  pdfauthor={Degree--diameter lab (Claude Code sessions coordinated by Doğa Özden)},
  hidelinks,
  pdfcreator={LaTeX via pandoc}}

\title{Bipartite graphs of diameter 7 from the split Cayley hexagon: level graphs, surgery and hubs}
\author{Degree--diameter lab (Claude Code sessions coordinated by Doğa Özden)}
\date{Draft, 25 October 2026}

\begin{document}
\maketitle

Draft of 25 October 2026, written for Ignacio López and Gabriela Araujo-Pardo by the degree--diameter lab
(github.com/dogaozden/degree-diameter; Claude Code sessions coordinated by Doğa Özden).

This note states and proves the constructions behind the lab's new lower bounds for the bipartite degree--diameter
problem at diameter 7 (cells (7,7), (9,7)--(16,7)). Everything below is meant to be checkable by hand: the geometric
hypotheses are stated in the language of the hexagon, and the neighbourhood of every vertex of every graph is given
explicitly (Definitions 2.1 and 3.1). Where a step depends on a computer search (the specific large sets and the
auxiliary ``hub graphs'' behind the record values), we say so and give the check. Section 10 lists the computer tests of
every theorem on small \(q\).

\emph{File paths} refer to the repository github.com/dogaozden/degree-diameter. Unless marked otherwise they are on the
master branch (commit c48c976); {[}b14{]} marks files on branch claude/cell-14-7-record-5d03gb and {[}b15{]} files on
branch claude/cell-15-7-haar-search-x6ynhl. This note and its tests (research/theorems/D7\_HEXAGON.md, d7hex\_tests/,
d7hex\_lemma41/) are on branch claude/d7-hexagon-theorems-4w1prt. Section 12 collects remarks for readers of the lab's
per-cell notes.

\hypertarget{contents}{%
\subsection{Contents}\label{contents}}

\begin{enumerate}
\def\labelenumi{\arabic{enumi}.}
\setcounter{enumi}{-1}
\tightlist
\item
  Results
\item
  The split Cayley hexagon: coordinates, distances, sets
\item
  Level graphs (classical)
\item
  Surgery and hubs: the general assembly theorem
\item
  Theorem C: transitions that extend to an automorphism (friendly and duality transitions)
\item
  Theorem B: local dualities; mirrored sets
\item
  Typed level layouts and slot counts
\item
  Hub graphs
\item
  The record graphs as corollaries
\item
  What is proved for every \(q\)
\item
  Computer checks
\item
  Prior art, and what is new
\item
  Appendix: notes for readers of the lab notes
\end{enumerate}

\begin{center}\rule{0.5\linewidth}{0.5pt}\end{center}

\hypertarget{results}{%
\subsection{0. Results}\label{results}}

Notation: \(n_b(\Delta, D)\) is the largest order of a bipartite graph with maximum degree \(\le \Delta\) and diameter
\(\le D\). For a prime power \(q\) put \[N = N(q) = \frac{q^6-1}{q-1} = (q+1)(q^4+q^2+1),\] the number of points (and of
lines) of a generalized hexagon of order \((q,q)\). In the 2013 table {[}MS13, Table 4{]} each cell below has the value
\((2r+1)\cdot 2N(q)\), with \(\Delta = q+1+r\) for the \((q,r)\) listed (the largest value of this form for that
\(\Delta\)); this is the order of the level graph of Section 2.

\begin{longtable}[]{@{}
  >{\raggedright\arraybackslash}p{(\columnwidth - 10\tabcolsep) * \real{0.0722}}
  >{\raggedright\arraybackslash}p{(\columnwidth - 10\tabcolsep) * \real{0.0825}}
  >{\raggedleft\arraybackslash}p{(\columnwidth - 10\tabcolsep) * \real{0.1031}}
  >{\raggedleft\arraybackslash}p{(\columnwidth - 10\tabcolsep) * \real{0.1134}}
  >{\raggedright\arraybackslash}p{(\columnwidth - 10\tabcolsep) * \real{0.2887}}
  >{\raggedright\arraybackslash}p{(\columnwidth - 10\tabcolsep) * \real{0.3402}}@{}}
\toprule\noalign{}
\begin{minipage}[b]{\linewidth}\raggedright
cell \((\Delta,7)\)
\end{minipage} & \begin{minipage}[b]{\linewidth}\raggedright
\(q\), \(r\)
\end{minipage} & \begin{minipage}[b]{\linewidth}\raggedleft
table (2013)
\end{minipage} & \begin{minipage}[b]{\linewidth}\raggedleft
ours
\end{minipage} & \begin{minipage}[b]{\linewidth}\raggedright
construction (this note)
\end{minipage} & \begin{minipage}[b]{\linewidth}\raggedright
parameters
\end{minipage} \\
\midrule\noalign{}
\endhead
\bottomrule\noalign{}
\endlastfoot
(7,7) & 5, 1 & 23,436 & \textbf{23,598} & 3 levels, one local duality (Cor. 8.1) & \(\lvert R\rvert = 81\) (72 pairwise
opposite + 9) \\
(9,7) & 7, 1 & 117,648 & \textbf{118,238} & same & \(\lvert R\rvert = 295\) (288 + 7) \\
(10,7) & 8, 1 & 224,694 & \textbf{225,152} & same & \(\lvert R\rvert = 229\) (217 + 12) \\
(11,7) & 9, 1 & 398,580 & \textbf{413,180} & 3 levels, duality, spread points freed (Cor. 8.2) &
\(\lvert F\rvert = 7{,}300\) \\
(12,7) & 9, 2 & 664,300 & \textbf{693,500} & 5 levels, duality, spread points freed (Cor. 8.2) &
\(\lvert F\rvert = 7{,}300\), \(14{,}600 + 14{,}600\) hubs \\
(13,7) & 11, 1 & 1,062,936 & \textbf{1,064,472} & 3 levels, one local duality (Cor. 8.1) & \(\lvert R\rvert = 768\) (756
+ 12) \\
(14,7) & 11, 2 & 1,771,560 & \textbf{1,789,588} & typed levels HDHDH + local duality (Cor. 8.3) & 8,645 pts / 8,246
lines / 768; \(9{,}014 + 9{,}014\) hubs \\
(15,7) & 11, 3 & 2,480,184 & \textbf{2,501,491} & typed levels HHDHDHD + local duality (Cor. 8.3) & twelve segment sets,
768; \(10{,}826 + 10{,}481\) hubs \\
(16,7) & 13, 2 & 4,022,340 & \textbf{4,052,108} & typed levels HDHDH + pairwise opposite sets (Cor. 8.3) & 14,762 pts /
14,823 lines / 122; \(14{,}884 + 14{,}884\) hubs \\
\end{longtable}

All nine graphs were rebuilt by the lab's coordinator from their solution files and checked by a breadth-first search
from every vertex (no symmetry assumed): bipartite, maximum degree \(\Delta\), diameter 7 (logs in
research/coordinator\_check/, Section 10). The cell (8,7) record is a Haar graph and is not discussed here.

Our comparison is with the 2013 table only. We could not read the current CombinatoricsWiki table from our environment
(the page returns HTTP 503), and a literature search found no later values for these cells; we would be grateful to hear
of any.

The theorems are:

\begin{itemize}
\tightlist
\item
  \textbf{Theorem A} (classical): the level graph with \(2r+1\) levels and arbitrary matchings has diameter \(\le 7\).
\item
  \textbf{Lemma 3.3 and Theorem 3.4} (assembly): if every ``transition'' survives the deletion of its freed matching
  edges with all its pairs still within distance 7, then hub vertices can be attached to the freed slots in any way, and
  the result has diameter at most 7 as soon as hubs of the same colour are within distance 6.
\item
  \textbf{Theorem C}: a transition whose matching is the restriction of an automorphism of the incidence graph (the
  identity between \(H\) and its dual, or a duality of \(H\)) survives the deletion at any \emph{segment set}: a point
  set in which every point is collinear with other points of the set along at most one line. Such sets have at most
  \((q+1)(q^3+1)\) points, with equality exactly for the point sets of spreads; and a transition of this kind is
  repaired only if its freed set is a segment set (Remark 4.2).
\item
  \textbf{Theorem B}: a transition with an arbitrary matching \(M\) survives the deletion at a non-collinear point set
  \(R\) if \(M(R)\) is non-concurrent and \(M\) maps the points collinear with \(x\) onto the lines concurrent with
  \(M(x)\), for every \(x\in R\). Propositions 5.1 and 5.2 say when such an \(M\) exists (``mirrored sets'').
\end{itemize}

\begin{center}\rule{0.5\linewidth}{0.5pt}\end{center}

\hypertarget{the-split-cayley-hexagon}{%
\subsection{1. The split Cayley hexagon}\label{the-split-cayley-hexagon}}

\hypertarget{coordinates}{%
\subsubsection{1.1 Coordinates}\label{coordinates}}

Let \(q\) be a prime power and \(V = \mathbb F_q^7\) with coordinates \((X_0,\dots,X_6)\). Let
\[Q(X) = X_0X_4 + X_1X_5 + X_2X_6 - X_3^2,\] with polar form
\(b(X,Y) = Q(X+Y)-Q(X)-Q(Y) = X_0Y_4+X_4Y_0+X_1Y_5+X_5Y_1+X_2Y_6+X_6Y_2-2X_3Y_3\) (alternating when \(q\) is even). The
parabolic quadric \(Q(6,q)\) is the set of points \(\langle x\rangle\) of \(\mathrm{PG}(6,q)\) with \(Q(x)=0\). It is
convenient to write \(x = (v, a, w)\) with \(v = (x_0,x_1,x_2)\), \(a = x_3\), \(w = (x_4,x_5,x_6)\); then
\(Q(x) = v\cdot w - a^2\).

The \textbf{split Cayley hexagon} \(H = H(q)\) ({[}Ti59{]}; see {[}VM98, §2.4{]}) has

\begin{itemize}
\tightlist
\item
  \textbf{points}: the set \(P\) of all \(N = (q^6-1)/(q-1)\) points of \(Q(6,q)\);
\item
  \textbf{lines}: the set \(L\) of the lines of \(Q(6,q)\) whose Grassmann coordinates \(p_{ij} = x_iy_j - x_jy_i\) (for
  any two spanning points \(x, y\)) satisfy Tits' equations, which in the labelling of coordinates used here read
  \[p_{12} = p_{34},\quad p_{54} = p_{32},\quad p_{20} = p_{35},\quad p_{65} = p_{30},\quad p_{01} = p_{36},\quad p_{46} = p_{31};\]
\item
  \textbf{incidence}: containment.
\end{itemize}

(Sources label the coordinates differently, and the equations change accordingly; this is the form used by all the lab's
generators. We checked by computer, from scratch, that with this quadric it gives a generalized hexagon of order
\((q,q)\) for \(q = 2, 3, 4, 5\): right numbers of points and lines, \(q+1\) on each, incidence graph of diameter 6 and
girth 12, Section 10.1.)

\emph{Example.} For \(x = e_0 = (1,0,\dots,0)\) and a point \(y\), \(p_{0j} = y_j\) and \(p_{ij} = 0\) for \(i,j\ge 1\);
the equations reduce to \(y_1 = y_2 = y_3 = 0\), and \(b(e_0,y) = y_4 = 0\). So the \(q+1\) lines through \(e_0\) are
\(\langle e_0, \lambda e_5+\mu e_6\rangle\), \((\lambda:\mu)\in\mathrm{PG}(1,q)\): they lie in the totally singular
plane \(\langle e_0, e_5, e_6\rangle\). (In general the lines of \(H(q)\) through a point \(x\) are the lines through
\(x\) of a totally singular plane \(\pi_x\subset x^\perp\).)

\(H(q)\) is a generalized hexagon of order \((q,q)\): \(N\) points and \(N\) lines, \(q+1\) points on a line and \(q+1\)
lines on a point. The group \(G_2(q)\) acts on it; for instance \(g\in SL(3,q)\) acts by
\((v,a,w)\mapsto (gv, a, g^{-T}w)\) (this is the symmetry used to check the (16,7) graph,
research/squads/cells/d16D7.md).

The \textbf{dual hexagon} \(H^D\) has the lines of \(H\) as points and the points of \(H\) as lines. A \textbf{duality}
of \(H\) is an isomorphism \(H\to H^D\), i.e.~a bijection \(\varphi\) of \(P\cup L\) with \(\varphi(P) = L\),
\(\varphi(L) = P\), preserving incidence. \(H(q)\) has a duality if and only if \(q = 3^h\) ({[}Ti59{]}; {[}VM98, Ch.
3{]}). For \(q = 9\) we use an explicit duality found by computer (research/d10D7/dual/dualfind.c: it extends a polarity
of the subhexagon \(H(3)\) over \(\mathbb F_3\subset\mathbb F_9\) by propagation along geodesics and then checks that
all \(N(q+1)\) incidences are preserved).

\hypertarget{the-incidence-graph-and-distances}{%
\subsubsection{1.2 The incidence graph and distances}\label{the-incidence-graph-and-distances}}

\(I = I(q)\) denotes the incidence graph of \(H\): vertex set \(P\cup L\) (\(2N\) vertices), \(x\sim y\) when \(x\) and
\(y\) are incident. \(d(x,y)\) is the distance in \(I\), and \(\Gamma_i(x) = \{y : d(x,y) = i\}\). Points are at even
distance from points, lines at odd distance from points. Two points are \emph{collinear} when \(d = 2\), two lines
\emph{concurrent} when \(d = 2\), and two elements are \emph{opposite} when \(d = 6\).

\(\Gamma_2(x)\) for a point \(x\) is the set of the \(q(q+1)\) points collinear with \(x\); for a line \(\ell\) it is
the set of the \(q(q+1)\) lines concurrent with \(\ell\). The sizes are \(\lvert\Gamma_i(x)\rvert = (q+1)q^{i-1}\) for
\(1\le i\le 5\) and \(\lvert\Gamma_6(x)\rvert = q^5\).

We use the following standard facts.

\textbf{(F1)} \(I\) is connected, bipartite, \((q+1)\)-regular, of diameter 6 and girth 12 (this is the definition of a
generalized hexagon of order \((q,q)\)).

\textbf{(F2) Geodesics.} If \(d(x,y)\le 5\) there is a unique geodesic from \(x\) to \(y\) (two would close a cycle of
length \(\le 10\)). If \(d(x,y) = 6\), then every \(z\in\Gamma_1(x)\) satisfies \(d(z,y) = 5\) (it is odd, and at least
\(6-1\)), so for every neighbour \(z\) of \(x\) there is exactly one geodesic from \(x\) to \(y\) through \(z\).

\textbf{(F3) Covering lemma.} Let \(x, y\) be elements of \(H\). There is \(z\in\Gamma_2(x)\) with \(d(z,y)\le 3\) if
\(x\) and \(y\) have different types, and \(d(z,y)\le 4\) if they have the same type. \emph{Proof.} If
\(2\le d(x,y)\le 6\), take for \(z\) the vertex at distance 2 from \(x\) on a geodesic from \(x\) to \(y\):
\(d(z,y) = d(x,y)-2\). If \(d(x,y) = 1\), take \(z\in\Gamma_1(y)\setminus\{x\}\) (then \(z\in\Gamma_2(x)\) and
\(d(z,y) = 1\)). If \(x = y\) any \(z\in\Gamma_2(x)\) has \(d(z,y) = 2\). The largest value is \(6-2 = 4\) for the same
type and \(5-2 = 3\) for different types. \(\square\)

\textbf{(F4) Orthogonality.} For points \(x, y\) of \(H(q)\): \(d(x,y)\le 4\) if and only if \(b(x,y) = 0\); that is,
the points not opposite \(x\) are the points of \(Q(6,q)\) in the tangent hyperplane \(x^\perp\) {[}VM98, §2.4{]}.
\emph{Proof.} If \(d(x,y)\le 4\), there is a point \(z\) with \(x\) and \(y\) on lines of \(H\) through \(z\) (take
\(z = x\) if \(d(x,y)\le 2\)); these lines lie in the totally singular plane \(\pi_z\) of §1.1, so \(b(x,y) = 0\).
Conversely, both sets have \(1+q+q^2+q^3+q^4\) points (\(x^\perp\cap Q(6,q)\) is a cone with vertex \(x\) over a
\(Q(4,q)\), and \(1 + q(q+1) + q^3(q+1)\) points are at distance \(\le 4\) from \(x\)), so the inclusion is an equality.
We re-checked (F4) by computer for all pairs of points for \(q = 2, 3, 4\) and for a sample of points for \(q = 5\)
(Section 10.1); the lab's large computations (research/d13D7/REPORT.md §3) also rely on it.

\hypertarget{sets-of-points-and-lines}{%
\subsubsection{1.3 Sets of points and lines}\label{sets-of-points-and-lines}}

\textbf{Definition 1.1.} A set of points of \(H\) is

\begin{itemize}
\tightlist
\item
  \emph{non-collinear} if no two of its points are collinear (a partial distance-2 ovoid);
\item
  a \emph{partial ovoid} if its points are pairwise opposite; an \emph{ovoid} if moreover it has \(q^3+1\) points.
\end{itemize}

Dually, a set of lines is \emph{non-concurrent}, a \emph{partial spread} (pairwise opposite lines), or a \emph{spread}
(\(q^3+1\) pairwise opposite lines).

By (F4), partial ovoids of \(H(q)\) are exactly the sets of pairwise non-orthogonal points of \(Q(6,q)\), i.e.~partial
ovoids of the polar space \(Q(6,q)\).

\textbf{Lemma 1.2 (standard examples).}

\begin{enumerate}
\def\labelenumi{(\alph{enumi})}
\tightlist
\item
  If \(\Sigma\) is a 3-space of \(\mathrm{PG}(6,q)\) with \(\Sigma\cap Q(6,q)\) an elliptic quadric \(Q^-(3,q)\), its
  \(q^2+1\) points form a partial ovoid of \(H(q)\), for every \(q\).
\item
  (Hermitian spread {[}Th80{]}.) If \(\Pi\) is a hyperplane with \(\Pi\cap Q(6,q)\) an elliptic quadric \(Q^-(5,q)\),
  the lines of \(H(q)\) contained in \(\Pi\) form a spread, for every \(q\).
\item
  If \(S\) is a spread, every line not in \(S\) meets exactly one line of \(S\).
\item
  If \(q = 3^h\) and \(\varphi\) is a duality, \(\varphi(S)\) is an ovoid for every spread \(S\).
\end{enumerate}

\emph{Proof.} (a) If \(x\ne y\) are points of \(Q^-(3,q)\) with \(b(x,y) = 0\), then
\(Q(\lambda x+\mu y) = \lambda\mu\, b(x,y) = 0\), so the line \(xy\) lies in \(Q^-(3,q)\), which contains no line. So
\(b(x,y)\ne 0\) and \(x, y\) are opposite by (F4).

\begin{enumerate}
\def\labelenumi{(\alph{enumi})}
\setcounter{enumi}{1}
\tightlist
\item
  is in {[}Th80{]}; we also checked it by computer for small \(q\) (Section 10.1), and the record builders re-check, for
  the sets they use, the hypotheses of the theorems they apply (Section 10.2).
\item
  A line meeting two lines of \(S\) would put them at distance \(\le 4\). So the lines of \(S\) together with the
  \(q(q+1)\) lines meeting each of them are \((q^3+1)(1+q(q+1)) = (q^3+1)(q^2+q+1) = N\) distinct lines, i.e.~all lines.
  (d) \(\varphi\) preserves distances. \(\square\)
\end{enumerate}

Ovoids of \(Q(6,q)\) (equivalently of \(H(q)\)) exist for \(q = 3^h\); they do not exist for \(q\) even {[}Th81{]} or
for \(q\) prime \(\ge 5\) {[}BGS06{]}. The largest partial ovoids are not known in general; the lab's record graphs use
partial ovoids found by computer (Section 8).

The following class of sets is the one that matters for Theorem C. It contains the non-collinear sets and the point sets
of partial spreads.

\textbf{Definition 1.3 (segment sets).} A set \(F\) of points is a \emph{segment set} if for every \(p\in F\) the points
of \(F\setminus\{p\}\) that are collinear with \(p\) all lie on one line through \(p\). Equivalently: there are no
\(p, p', p''\in F\) with \(d(p,p') = d(p',p'') = 2\) and \(d(p,p'') = 4\) (a path of two collinearities that turns at
\(p'\)). Dually, a set \(G\) of lines is a \emph{dual segment set} if for every \(\ell\in G\) the lines of
\(G\setminus\{\ell\}\) that meet \(\ell\) all pass through one point of \(\ell\).

So a segment set is a disjoint union of ``segments'' (sets of \(\ge 2\) points on a line) and isolated points, such that
no two points from different pieces are collinear. Examples:

\begin{itemize}
\tightlist
\item
  every non-collinear set;
\item
  every set in which each point is collinear with at most one other point of the set (the ``relaxed sets'' (R1) of the
  lab notes);
\item
  every subset of the point set of a partial spread (points on distinct lines of a partial spread are not collinear,
  since the lines are opposite);
\item
  dually: the set of all lines through the points of a partial ovoid (lines through distinct points of a partial ovoid
  never meet).
\end{itemize}

\textbf{Proposition 1.4 (size of segment sets).} A segment set \(F\) has \(\lvert F\rvert\le (q+1)(q^3+1)\), with
equality if and only if \(F\) is the point set of a spread. Dually, a dual segment set has at most \((q+1)(q^3+1)\)
lines, with equality if and only if it is the set of lines through the points of an ovoid (so equality is possible only
when \(H(q)\) has an ovoid, e.g.~\(q = 3^h\), and impossible for \(q\) even and for \(q\) prime \(\ge 5\)).

\emph{Proof.} Let \(F_1\) be the points of \(F\) collinear with no other point of \(F\), and \(C_1,\dots,C_m\) the
segments (maximal sets of \(\ge 2\) points of \(F\) on a line \(\ell_i\); \(\lvert C_i\rvert\le q+1\)). A line meeting
\(F\) in exactly one point is called a \emph{private} line of that point. A point of \(F_1\) has \(q+1\) private lines
and a point of a segment \(C_i\) has \(q\) (all its lines except \(\ell_i\): by the segment property no other line
through it contains a second point of \(F\)). Private lines of different points are different, and different from the
\(\ell_i\). So \[(q+1)\lvert F_1\rvert + q\sum_i \lvert C_i\rvert + m\ \le\ N .\] Since
\(m\ge \sum_i\lvert C_i\rvert/(q+1)\), the left side is at least \((q + \tfrac1{q+1})\lvert F\rvert\), which gives
\(\lvert F\rvert\le N(q+1)/(q^2+q+1) = (q+1)(q^3+1)\). Equality forces \(F_1 = \emptyset\), \(\lvert C_i\rvert = q+1\)
(whole lines) and \(m = q^3+1\). The lines \(\ell_i\) are then pairwise opposite: if \(\ell_i\) and \(\ell_k\) met in a
point \(z\), then \(z\in F\) would have other points of \(F\) on two lines; if \(d(\ell_i,\ell_k) = 4\), a point of
\(\ell_i\) would be collinear (along the line joining them) with a point of \(\ell_k\), and would also have other points
of \(F\) on \(\ell_i\). So they form a spread. Conversely the point set of a spread is a segment set of this size. The
dual statement is the same proof in \(H^D\). \(\square\)

For comparison, a non-collinear set has at most \(N/(q+1) = q^4+q^2+1\) points (each line contains at most one of its
points; equality means a distance-2 ovoid, a point set meeting every line exactly once), and
\((q+1)(q^3+1) = q^4+q^3+q+1\) is larger by \(q^3-q^2+q\).

\begin{center}\rule{0.5\linewidth}{0.5pt}\end{center}

\hypertarget{level-graphs}{%
\subsection{2. Level graphs}\label{level-graphs}}

Fix \(r\ge 1\); there are \(2r+1\) levels, indexed by \(j\in\mathbb Z_{2r+1}\). Each level gets a \textbf{type}
\(\tau_j\in\{H, D\}\). For a type \(\tau\) let \[A(H) = P,\quad B(H) = L,\qquad A(D) = L,\quad B(D) = P\] (the
``points'' and ``lines'' of \(H\) or of \(H^D\); the letters \(H\), \(D\) for the types refer to the hexagon and its
dual, not to the diameter). A \textbf{transition} is a pair \(t = (j,k)\) with \(j\in\mathbb Z_{2r+1}\),
\(1\le k\le r\); its source level is \(j\) and its target level is \(j+k\). It is \textbf{friendly} if
\(\tau_j\ne\tau_{j+k}\); then \(A(\tau_j) = B(\tau_{j+k})\).

\textbf{Definition 2.1 (level graph).} Given types \(\tau\) and, for every transition \(t=(j,k)\), a bijection
\(M_t: A(\tau_j)\to B(\tau_{j+k})\), the graph \(X = X(\tau, M)\) has vertex set \((P\cup L)\times\mathbb Z_{2r+1}\) and
edges

\begin{itemize}
\tightlist
\item
  \((x,j)\sim(y,j)\) whenever \(x\) and \(y\) are incident (``level edges'');
\item
  \((a,j)\sim(M_{(j,k)}(a), j+k)\) for every \(j\), every \(1\le k\le r\) and every \(a\in A(\tau_j)\) (``matching
  edges'').
\end{itemize}

Colour the vertices \((a,j)\), \(a\in A(\tau_j)\), with colour \(\mathcal A\) and the vertices \((b,j)\),
\(b\in B(\tau_j)\), with colour \(\mathcal B\). Level edges and matching edges join different colours, so \(X\) is
bipartite. The neighbourhoods are:
\[N_X(a,j) = \Gamma_1(a)\times\{j\}\ \cup\ \{(M_{(j,k)}(a), j+k) : 1\le k\le r\}\qquad (a\in A(\tau_j)),\]
\[N_X(b,j) = \Gamma_1(b)\times\{j\}\ \cup\ \{(M_{(j-k,k)}^{-1}(b), j-k) : 1\le k\le r\}\qquad (b\in B(\tau_j)).\] So
\(X\) is \((q+1+r)\)-regular of order \((2r+1)\cdot 2N\).

When all levels have type \(H\) and all \(M_t\) equal one bijection \(\sigma: P\to L\), this is the graph whose order is
the table value (for \(r=1\): three copies of \(I\) joined by \((p,i)\sim(\sigma(p),i+1)\)).

\textbf{Theorem A (level graph).} For all types and all bijections \(M_t\), \(X(\tau, M)\) has diameter at most 7. More
precisely: two vertices of one level are joined by a path of length \(\le 6\) inside that level, and for every
transition \(t=(j,k)\) every vertex of level \(j\) is joined to every vertex of level \(j+k\) by a path of length
\(\le 7\) that uses only level edges of levels \(j\) and \(j+k\) and matching edges of \(t\).

\emph{Proof.} Inside a level this is (F1). Every two distinct levels differ by \(\pm k\) for some \(1\le k\le r\) (as
there are \(2r+1\) levels), so it is enough to join \((u,j)\) to \((v,j')\), \(j' = j+k\), with \(M = M_{(j,k)}\):

\begin{enumerate}
\def\labelenumi{\arabic{enumi}.}
\tightlist
\item
  \(u\in A(\tau_j)\): \((u,j)\to(M(u),j')\leadsto(v,j')\), length \(\le 1+6\).
\item
  \(u\in B(\tau_j)\), \(v\in B(\tau_{j'})\): with \(a = M^{-1}(v)\), \((u,j)\leadsto(a,j)\to(v,j')\), length \(\le 5+1\)
  (in \(I\), elements of different types are at distance \(\le 5\)).
\item
  \(u\in B(\tau_j)\), \(v\in A(\tau_{j'})\): pick \(b\in\Gamma_1(v)\) and \(a = M^{-1}(b)\):
  \((u,j)\leadsto(a,j)\to(b,j')\to(v,j')\), length \(\le 5+1+1\). \(\square\)
\end{enumerate}

\begin{center}\rule{0.5\linewidth}{0.5pt}\end{center}

\hypertarget{surgery-and-hubs}{%
\subsection{3. Surgery and hubs}\label{surgery-and-hubs}}

\hypertarget{the-construction}{%
\subsubsection{3.1 The construction}\label{the-construction}}

\textbf{Definition 3.1 (surgery and hubs).} Start from a level graph \(X(\tau, M)\) and choose for every transition
\(t=(j,k)\) a set \(R_t\subseteq A(\tau_j)\) of \textbf{freed elements} (possibly empty).

\begin{enumerate}
\def\labelenumi{\arabic{enumi}.}
\tightlist
\item
  \emph{Deletion.} Delete the matching edges \((a,j)\sim(M_t(a),j+k)\), \(a\in R_t\). Call the result \(X_R\). Each
  deleted edge leaves a \textbf{free \(\mathcal A\)-slot} \((a,t)\) at the vertex \((a,j)\) and a \textbf{free
  \(\mathcal B\)-slot} \((M_t(a),t)\) at the vertex \((M_t(a), j+k)\). (An element freed on several transitions out of,
  or into, a level carries several slots there.) Level \(j\) has
  \[\alpha_j = \sum_{k=1}^r \lvert R_{(j,k)}\rvert\ \text{ free }\mathcal A\text{-slots},\qquad
    \beta_j = \sum_{k=1}^r \lvert R_{(j-k,k)}\rvert\ \text{ free }\mathcal B\text{-slots}.\]
\item
  \emph{Hubs.} Add \(K_{\mathcal B}\le\min_j\alpha_j\) new vertices of colour \(\mathcal B\)
  (\textbf{\(\mathcal B\)-hubs}) and \(K_{\mathcal A}\le\min_j\beta_j\) new vertices of colour \(\mathcal A\)
  (\textbf{\(\mathcal A\)-hubs}). Each \(\mathcal B\)-hub \(h\) takes exactly one free \(\mathcal A\)-slot in every
  level \(j\) and is joined to its vertex, which we write \((\omega_j(h), j)\); each \(\mathcal A\)-hub \(g\) likewise
  takes one free \(\mathcal B\)-slot in every level, at the vertex \((\omega'_j(g), j)\). Every slot is taken at most
  once. (The assignment is otherwise arbitrary; a hub may take slots of different freed elements in different levels,
  and a vertex with several free slots may receive several hubs.)
\item
  \emph{Hub graph.} Add a bipartite graph \(Y\) between the \(\mathcal B\)-hubs and the \(\mathcal A\)-hubs.
\end{enumerate}

Call the result \(X_R^+\). In the lab notes, \(\mathcal B\)-hubs are ``line-side hubs'' (in an \(H\) level they are
joined to points) and \(\mathcal A\)-hubs ``point-side hubs''.

\emph{Exact neighbourhoods.} For a level vertex, the neighbourhood is that of Definition 2.1 minus the deleted matching
edges plus one hub for each taken free slot. For hubs:
\[N_{X_R^+}(h) = \{(\omega_j(h), j) : j\in\mathbb Z_{2r+1}\}\cup N_Y(h),\qquad N_{X_R^+}(g) = \{(\omega'_j(g), j): j\in\mathbb Z_{2r+1}\}\cup N_Y(g).\]

\emph{Degree and order.} Level vertices have degree \(\le q+1+r\) (equality when all their free slots are used), hubs
have degree \(2r+1+\deg_Y\). So \(X_R^+\) is bipartite with maximum degree \(\le q+1+r\) as soon as \(Y\) has maximum
degree \(\le q-r\), and its order is \[\lvert X_R^+\rvert = (2r+1)\cdot 2N + K_{\mathcal A} + K_{\mathcal B}.\]

\hypertarget{assembly}{%
\subsubsection{3.2 Assembly}\label{assembly}}

\textbf{Definition 3.2.} A transition \(t=(j,k)\) is \textbf{repaired} (by \(R_t\)) if in \(X_R\) every vertex of level
\(j\) is joined to every vertex of level \(j+k\) by a path of length \(\le 7\) that uses only level edges of levels
\(j, j+k\) and the remaining matching edges of \(t\).

Whether \(t\) is repaired depends only on \(M_t\) and \(R_t\) (not on the other transitions, and not on the hubs).

\textbf{Lemma 3.3 (hubs).} Let \(X_R^+\) be as in Definition 3.1. Then

\begin{itemize}
\tightlist
\item
  every hub is within distance 7 of every level vertex;
\item
  every \(\mathcal B\)-hub is within distance 7 of every \(\mathcal A\)-hub.
\end{itemize}

\emph{Proof.} A \(\mathcal B\)-hub \(h\) and a level vertex \((v,j)\): \(h\to(\omega_j(h),j)\leadsto(v,j)\) inside level
\(j\), length \(\le 1+6\). The level copies are never touched by the surgery. A \(\mathcal B\)-hub \(h\) and an
\(\mathcal A\)-hub \(g\): in level \(0\), \(\omega_0(h)\in A(\tau_0)\) and \(\omega'_0(g)\in B(\tau_0)\) have different
types, so \(h\to(\omega_0(h),0)\leadsto(\omega'_0(g),0)\to g\) has length \(\le 1+5+1\). \(\square\)

\textbf{Theorem 3.4 (assembly).} Suppose

\begin{itemize}
\item
  \begin{enumerate}
  \def\labelenumi{(\Alph{enumi})}
  \setcounter{enumi}{19}
  \tightlist
  \item
    every transition \(t\) is repaired by \(R_t\), and
  \end{enumerate}
\item
  (HG) any two \(\mathcal B\)-hubs, and any two \(\mathcal A\)-hubs, are within distance 6 in \(X_R^+\).
\end{itemize}

Then \(X_R^+\) has diameter \(\le 7\). If moreover \(Y\) has maximum degree \(\le q-r\), then
\[n_b(q+1+r,\,7)\ \ge\ (2r+1)\cdot 2N(q) + K_{\mathcal A} + K_{\mathcal B}.\]

\emph{Proof.} Pairs of level vertices: same level by (F1), different levels by (T) (as in the proof of Theorem A, every
pair of levels is the pair of some transition). Pairs with a hub: Lemma 3.3 and (HG). \(\square\)

(HG) holds for instance when \(Y\) alone has all same-colour distances \(\le 6\); it can also use paths through the
levels (two \(\mathcal B\)-hubs whose slots in some level are at distance \(\le 4\) are within distance 6). Section 7
discusses hub graphs.

So everything reduces to geometric conditions under which one transition is repaired. Sections 4 and 5 give two.

\begin{center}\rule{0.5\linewidth}{0.5pt}\end{center}

\hypertarget{theorem-c-transitions-that-extend-to-an-automorphism}{%
\subsection{4. Theorem C: transitions that extend to an
automorphism}\label{theorem-c-transitions-that-extend-to-an-automorphism}}

A transition \(t = (j,k)\) is \textbf{automorphic} if \(M_t\) is the restriction to \(A(\tau_j)\) of an automorphism
\(\psi\) of the graph \(I\) with \(\psi(A(\tau_j)) = B(\tau_{j+k})\). The two cases used here:

\begin{itemize}
\tightlist
\item
  a \emph{friendly} transition with \(M_t = \mathrm{id}\) (\(\psi = \mathrm{id}\); \(A(\tau_j) = B(\tau_{j+k})\));
\item
  an \(H\to H\) transition with \(M_t = \varphi\vert_P\) for a duality \(\varphi\) (\(q = 3^h\)); a \(D\to D\)
  transition with \(M_t = \varphi\vert_L\).
\end{itemize}

(A friendly transition whose matching is a collineation of \(H\), restricted to points or to lines, is also
automorphic.)

\textbf{Lemma 4.1.} Let \(F\) be a segment set of points. For all elements \(u, v\) of \(H\) there is a point
\(p\notin F\) with \[d(u,p) + d(p,v)\ \le\ 6 .\]

\emph{Proof.} If some point \(p\) on a geodesic from \(u\) to \(v\) is not in \(F\), then
\(d(u,p)+d(p,v) = d(u,v)\le 6\). Assume that all points on all geodesics from \(u\) to \(v\) lie in \(F\).

A geodesic with three points \(p, p', p''\) in a row (\(p - m - p' - m' - p''\) with lines \(m\ne m'\)) cannot have all
three in \(F\) (that is a turning path, excluded by Definition 1.3). A geodesic between two points at distance 4 or 6, a
point and a line at distance 5, or two lines at distance 6 contains three points in a row. So only the following cases
remain; in each we exhibit \(p\notin F\) with \(d(u,p)+d(p,v)\le 6\). We use repeatedly: \emph{if \(p_0\in F\) has a
point of \(F\) on the line \(s\) through it, then no point \(p\ne p_0\) collinear with \(p_0\) off \(s\) is in \(F\)}
(Definition 1.3).

\begin{enumerate}
\def\labelenumi{(\alph{enumi})}
\item
  \(u = v\) a point of \(F\). Let \(s\) be a line through \(u\) containing every point of \(F\) collinear with \(u\)
  (any line if there is none). Take a line \(y\ne s\) through \(u\) and \(p\in y\setminus\{u\}\): \(p\notin F\) and
  \(d(u,p)+d(p,u) = 4\).
\item
  \(u, v\) collinear points of \(F\) on the line \(s\). Take \(y\ne s\) through \(u\) and \(p\in y\setminus\{u\}\):
  \(p\notin F\), \(d(u,p) = 2\), \(d(p,v)\le 4\).
\item
  \(u\in F\) a point, \(v\) a line through \(u\). If \(v\) has a point \(p\ne u\) outside \(F\): \(2+1 = 3\). Otherwise
  all points of \(v\) are in \(F\); take \(y\ne v\) through \(u\) and \(p\in y\setminus\{u\}\): \(p\notin F\),
  \(d(u,p) = 2\), \(d(p,v)\le 3\).
\item
  \(u\) a point, \(v\) a line, \(d(u,v) = 3\), geodesic \(u - s - p_2 - v\) with \(u, p_2\in F\). A point
  \(p\in v\setminus\{p_2\}\) has \(d(u,p) = 4\), \(d(p,v) = 1\). If all of them were in \(F\), \(p_2\) would have points
  of \(F\) on the two lines \(s\) and \(v\).
\item
  \(u, v\) lines meeting in \(p_1\in F\). Not both \(u\) and \(v\) contain a second point of \(F\) (else \(p_1\) turns).
  Say \(u\) does not: take \(p\in u\setminus\{p_1\}\), so \(p\notin F\), \(d(u,p) = 1\), \(d(p,v) = 3\).
\item
  \(u, v\) lines at distance 4, geodesic \(u - p_1 - y - p_3 - v\) with \(p_1, p_3\in F\). Then no other point \(p\) of
  \(u\) is in \(F\) (else \(p_1\) turns: \(p\) on \(u\), \(p_3\) on \(y\)), and \(d(u,p) = 1\),
  \(d(p,v)\le d(p,p_1)+d(p_1,v) = 2+3\).
\item
  \(u = v\) a line. If some point of \(u\) is not in \(F\): \(1+1\). Otherwise take \(p_1\in u\), a line \(y\ne u\)
  through \(p_1\) and \(p\in y\setminus\{p_1\}\): \(p\notin F\) and \(d(u,p) = 3\), so \(3+3 = 6\).
\end{enumerate}

The cases with \(u\) a line and \(v\) a point are (c), (d) with \(u\) and \(v\) exchanged. \(\square\)

\textbf{Theorem C.} Let \(t = (j,k)\) be an automorphic transition. If \(A(\tau_j) = P\) and \(R_t\) is a segment set of
points, or \(A(\tau_j) = L\) and \(R_t\) is a dual segment set of lines, then \(t\) is repaired.

\emph{Proof.} Let \(\psi\) be the automorphism of \(I\) extending \(M = M_t\), \(j' = j+k\), and take \((u,j)\),
\((v,j')\). For every \(a\in A(\tau_j)\setminus R_t\) the path \[(u,j)\leadsto(a,j)\to(\psi(a),j')\leadsto(v,j')\] has
length \(d(u,a) + 1 + d(\psi(a),v) = d(u,a) + 1 + d(a,\psi^{-1}(v))\), since \(\psi\) is an isometry of \(I\). If
\(A(\tau_j) = P\), Lemma 4.1 applied to \(u\), \(\psi^{-1}(v)\) and \(F = R_t\) gives \(a\) with
\(d(u,a)+d(a,\psi^{-1}(v))\le 6\), so the path has length \(\le 7\). If \(A(\tau_j) = L\), apply Lemma 4.1 in \(H^D\)
(same graph \(I\), points and lines exchanged; a dual segment set of \(H\) is a segment set of \(H^D\)). \(\square\)

\textbf{Remark 4.2 (the converse).} An automorphic transition is repaired (Definition 3.2) \emph{if and only if} its
freed set is a segment set (dual segment set for lines). The ``if'' is Theorem C. For ``only if'', suppose that
\(p, p', p''\in R_t\) turn at \(p'\) (\(p - m - p' - m' - p''\), \(m\ne m'\)), take a line \(u\ne m'\) through \(p''\).
Then \(d(p,u) = 5\) with the unique geodesic \(p - m - p' - m' - p'' - u\). A path from \((u,j)\) to \((\psi(p),j')\)
inside the two levels uses an odd number \(2h+1\) of matching edges of \(t\), and (transporting the segments in level
\(j'\) back by \(\psi^{-1}\) and using the triangle inequality) has length at least \(2h+1+d(u,p) = 2h+6\). So a path of
length \(\le 7\) uses one matching edge, at a point \(a\) with \(d(u,a)+d(a,p) = 5\), i.e.~\(a\in\{p, p', p''\}\): all
three edges are deleted, and the distance in the two levels is \(> 7\). (In the whole graph, routes through a third
level or through hubs might still exist; in the lab's five-level graph over \(H(9)\) with the 7,300 spread points freed
and no hubs, freeing one more point left 82,720 ordered pairs at distance \(> 7\), research/squads/cells/d12D7.md §2.)

\textbf{Corollary 4.3.} The following transitions are repaired:

\begin{enumerate}
\def\labelenumi{\arabic{enumi}.}
\tightlist
\item
  a friendly transition with \(M_t = \mathrm{id}\) and \(R_t\) a non-collinear set of points (if \(\tau_j = H\)) or a
  non-concurrent set of lines (if \(\tau_j = D\)) --- the ``one-transition theorem at friendly transitions'' of the lab
  notes;
\item
  the same with \(R_t\) a set in which every element has at most one partner (collinear, resp. concurrent element of
  \(R_t\)) --- the relaxed sets (R1) of research/d15D7b/RELAXED.md, \textbf{without} its hub condition (R2);
\item
  an \(H\to H\) transition with \(M_t\) a duality and \(R_t\) any subset of the points of a partial spread
  (research/squads/cells/d12D7.md §2; e.g.~all \((q+1)(q^3+1)\) points of the Hermitian spread, \(q=3^h\)).
\end{enumerate}

\begin{center}\rule{0.5\linewidth}{0.5pt}\end{center}

\hypertarget{theorem-b-local-dualities}{%
\subsection{5. Theorem B: local dualities}\label{theorem-b-local-dualities}}

For transitions that are not automorphic (an \(H\to H\) transition when \(q\ne 3^h\), or any \(H\to H\) transition whose
matching is not a duality) we only ask \(M_t\) to behave like a duality near the freed elements.

\textbf{Theorem B (one-transition surgery).} Let \(t = (j,k)\) be any transition, \(M = M_t\), and
\(R = R_t\subseteq A(\tau_j)\). Suppose

\begin{itemize}
\tightlist
\item
  (H1) no two elements of \(R\) are at distance 2 (non-collinear, resp. non-concurrent);
\item
  (H2) no two elements of \(M(R)\) are at distance 2;
\item
  (H3) \(M(\Gamma_2(x)) = \Gamma_2(M(x))\) for every \(x\in R\) (\emph{\(M\) is a local duality at \(x\)} when \(M\)
  maps points to lines: it maps the points collinear with \(x\) onto the lines concurrent with \(M(x)\)).
\end{itemize}

Then \(t\) is repaired.

\emph{Proof.} Write \(j' = j+k\); ``\(\mathrm A\)'' and ``\(\mathrm B\)'' refer to \(A(\tau_j), B(\tau_j)\) in level
\(j\) and to \(A(\tau_{j'}), B(\tau_{j'})\) in level \(j'\). Elements of different types are at distance \(\le 5\) in
\(I\). Take \((u,j)\) and \((v,j')\).

\begin{enumerate}
\def\labelenumi{\arabic{enumi}.}
\tightlist
\item
  \(u\in \mathrm A\), \(u\notin R\): \((u,j)\to(M(u),j')\leadsto(v,j')\): \(\le 1+6\).
\item
  \(u\in\mathrm B\), \(v\in\mathrm A\): the \(q+1\ge 2\) neighbours of \(v\) are pairwise at distance 2, so by (H2) at
  most one is in \(M(R)\); take \(b\in\Gamma_1(v)\setminus M(R)\) and \(a = M^{-1}(b)\notin R\):
  \((u,j)\leadsto(a,j)\to(b,j')\to(v,j')\): \(\le 5+1+1\).
\item
  \(u\in\mathrm B\), \(v\in\mathrm B\setminus M(R)\): \((u,j)\leadsto(M^{-1}(v),j)\to(v,j')\): \(\le 5+1\).
\item
  \(u\in\mathrm B\), \(v = M(x)\), \(x\in R\): by (F3) there is \(z\in\Gamma_2(x)\) with \(d(z,u)\le 3\) (\(x\) and
  \(u\) have different types). By (H1) \(z\notin R\); by (H3) \(M(z)\in\Gamma_2(v)\):
  \((v,j')\leadsto(M(z),j')\to(z,j)\leadsto(u,j)\): \(\le 2+1+3\).
\item
  \(u = x\in R\), \(v\in\mathrm A\): by (F3) there is \(w\in\Gamma_2(M(x))\) with \(d(w,v)\le 3\); by (H3) \(w = M(z)\)
  with \(z\in\Gamma_2(x)\), and \(z\notin R\) by (H1): \((x,j)\leadsto(z,j)\to(w,j')\leadsto(v,j')\): \(\le 2+1+3\).
\item
  \(u = x\in R\), \(v\in\mathrm B\): the same path, now with \(d(w,v)\le 4\) by (F3) (\(M(x)\) and \(v\) have the same
  type): \(\le 2+1+4\).
\end{enumerate}

These cases cover all pairs (\(u\in\mathrm A\setminus R\): 1; \(u\in R\): 5, 6; \(u\in\mathrm B\): 2, 3, 4). \(\square\)

For a friendly transition with \(M = \mathrm{id}\), (H3) holds at every element and (H2) is (H1), so Theorem B gives
Corollary 4.3.1 again; Theorem C is stronger there. Theorem B is needed for non-automorphic transitions. We now describe
its hypotheses for an \(H\to H\) transition (\(M: P\to L\)); the \(D\to D\) case is the same with \(M^{-1}\).

For two points \(x, y\) at distance 4 let \(\mathrm{mid}(x,y)\) be the unique point collinear with both; for two lines
at distance 4 let \(\mathrm{mid}(\ell,m)\) be the unique line concurrent with both.

\textbf{Proposition 5.1 (what (H1)--(H3) force).} Under (H1)--(H3), \(M\) restricted to \(R\) preserves distances, and
\(M(\mathrm{mid}(x,y)) = \mathrm{mid}(M(x), M(y))\) for all \(x, y\in R\) at distance 4.

\emph{Proof.} Let \(x\ne y\in R\); \(d(x,y)\in\{4,6\}\) by (H1). If \(d(x,y) = 4\),
\(m = \mathrm{mid}(x,y)\in\Gamma_2(x)\cap\Gamma_2(y)\), so by (H3) \(M(m)\in\Gamma_2(M(x))\cap\Gamma_2(M(y))\); hence
\(d(M(x),M(y))\le 4\), \(=4\) by (H2), and \(M(m)\) is the mid-line. If \(d(x,y) = 6\),
\(\Gamma_2(x)\cap\Gamma_2(y) = \emptyset\), so \(\Gamma_2(M(x))\cap\Gamma_2(M(y)) = \emptyset\) by (H3), and
\(d(M(x),M(y))\notin\{2,4\}\), i.e.~\(= 6\). \(\square\)

So \(M(R)\) must \emph{mirror} the distance-4 structure of \(R\). The converse, in the form used for the records:

\textbf{Proposition 5.2 (mirrored sets).} Let \(R\) be a non-collinear set of points, \(R^*\) a non-concurrent set of
lines and \(f: R\to R^*\) a bijection such that

\begin{itemize}
\tightlist
\item
  (S1) every point is collinear with at most two points of \(R\), and every line is concurrent with at most two lines of
  \(R^*\);
\item
  (S2) for \(x\ne y\in R\): \(d(x,y) = 4\) if and only if \(d(f(x),f(y)) = 4\).
\end{itemize}

Then there is a bijection \(M: P\to L\) with \(M\vert_R = f\) satisfying (H1)--(H3) for \(R\). One such \(M\) is:
\(M(x) = f(x)\) for \(x\in R\); \(M(\mathrm{mid}(x,y)) = \mathrm{mid}(f(x),f(y))\) for \(x,y\in R\) at distance 4; for
each \(x\in R\), any bijection from the remaining points of \(\Gamma_2(x)\) onto the remaining lines of
\(\Gamma_2(f(x))\); any bijection on the rest.

\emph{Proof.} Let \(U = \bigcup_{x\in R}\Gamma_2(x)\) and \(U^* = \bigcup_{x\in R}\Gamma_2(f(x))\). By (H1) for \(R\)
and \(R^*\), \(R\cap U = \emptyset\) and \(R^*\cap U^* = \emptyset\). A point of \(U\) lies in \(\Gamma_2(x)\) for one
or two \(x\in R\) (S1); if two, \(x\) and \(y\) are at distance 4 (not collinear, and both collinear with it) and the
point is \(\mathrm{mid}(x,y)\). So the points of \(U\) in two sets \(\Gamma_2(\cdot)\) correspond bijectively to the
pairs \(\{x,y\}\subseteq R\) at distance 4, and likewise for \(U^*\); by (S2) \(\{x,y\}\mapsto\{f(x),f(y)\}\) is a
bijection between these pairs. Hence the rule \(\mathrm{mid}(x,y)\mapsto\mathrm{mid}(f(x),f(y))\) is a bijection from
the shared points of \(U\) to the shared lines of \(U^*\), and it maps the shared points of \(\Gamma_2(x)\) onto the
shared lines of \(\Gamma_2(f(x))\). The remaining (``private'') points of \(\Gamma_2(x)\) and private lines of
\(\Gamma_2(f(x))\) are then equally many (both \(q(q+1)\) minus the number of \(y\in R\) at distance 4 from \(x\)), and
private sets of different \(x\) are disjoint. So the pieces define a bijection \(R\cup U\to R^*\cup U^*\) with
\(M(\Gamma_2(x)) = \Gamma_2(f(x))\) for every \(x\in R\). As
\(\lvert P\setminus (R\cup U)\rvert = \lvert L\setminus(R^*\cup U^*)\rvert\), it extends to a bijection \(P\to L\).
(H1), (H2) are the hypotheses on \(R\) and \(R^*\). \(\square\)

\textbf{Corollary 5.3.}

\begin{enumerate}
\def\labelenumi{(\alph{enumi})}
\tightlist
\item
  (Pairwise opposite sets; research/d13D7/REPORT.md.) If \(R\) is a partial ovoid and \(R^*\) a partial spread with
  \(\lvert R\rvert = \lvert R^*\rvert\), every bijection \(f\) satisfies (S1), (S2) (there are no pairs at distance 4).
\item
  (Mirrored sets; research/cell147/mirror/.) Let \(R_0\) and \(X\) be disjoint partial ovoids, \(f_0: R_0\to S_0\) a
  bijection onto a partial spread \(S_0\), and \(x\mapsto L_x\) an injective map from \(X\) onto a partial spread
  disjoint from \(S_0\). Put \(R = R_0\cup X\), \(R^* = S_0\cup\{L_x : x\in X\}\), \(f = f_0\) on \(R_0\) and
  \(f(x) = L_x\) on \(X\), and assume \(R\) non-collinear and \(R^*\) non-concurrent. Then (S1) holds automatically:
  three points of \(R\) collinear with one point would be pairwise at distance 4, but two of them lie in \(R_0\) or two
  in \(X\), which are pairwise opposite; likewise for lines. And (S2) says: for \(r\in R_0\) and \(x\in X\),
  \(d(r,x) = 4\) if and only if \(d(f_0(r), L_x) = 4\).
\end{enumerate}

In (b) the matching condition is a bipartite assignment problem: each \(r\in R_0\) has a ``type'', the set of \(x\in X\)
at distance 4 from it, and must be sent to a spread line of the same type with respect to the lines \(L_x\). The record
sets were found by annealing over \(X\), \(L_x\) and \(f\) (research/cell147/mirror/mirror\_sa*.py {[}b14{]}), and every
one was re-checked by the builder (non-collinearity, non-concurrency, and (H3) at every \(x\in R\), directly).

\begin{center}\rule{0.5\linewidth}{0.5pt}\end{center}

\hypertarget{typed-level-layouts-and-slot-counts}{%
\subsection{6. Typed level layouts and slot counts}\label{typed-level-layouts-and-slot-counts}}

By Theorem 3.4 the order is \((2r+1)\cdot 2N + K_{\mathcal A} + K_{\mathcal B}\) with
\(K_{\mathcal B}\le\min_j\alpha_j\) and \(K_{\mathcal A}\le\min_j\beta_j\), where \(\alpha_j\) (\(\beta_j\)) counts the
freed elements of the transitions leaving (entering) level \(j\). Which transitions are friendly depends on the types.

\textbf{All levels of type \(H\), one bijection \(\sigma\), one freed set \(R\) on every transition} (the three-level
records (7,7), (9,7), (10,7), (13,7) with \(\sigma\) a local duality on \(R\), and the duality records (11,7) with three
levels and (12,7) with five). Every level has \(\alpha_j = \beta_j = r\lvert R\rvert\), so
\(K_{\mathcal A} = K_{\mathcal B} = r\lvert R\rvert\) and \[\lvert X_R^+\rvert = (2r+1)\cdot 2N + 2r\lvert R\rvert .\]
Hubs: for \(r = 1\) the \(\mathcal B\)-hub \(h_x\) of \(x\in R\) is joined to \((x,0),(x,1),(x,2)\) and the
\(\mathcal A\)-hub \(g_x\) to \((\sigma(x),0),(\sigma(x),1),(\sigma(x),2)\), plus \(q-1\) hub-graph edges each.

\textbf{Types \(HDHDH\) (\(r = 2\)).} Transitions \(j\to j+1\) and \(j\to j+2\). Friendly: \(0\to1\), \(2\to3\),
\(4\to1\) (\(H\to D\), points freed) and \(1\to2\), \(3\to4\), \(3\to0\) (\(D\to H\), lines freed). Not friendly:
\(0\to2\), \(2\to4\), \(4\to0\) (\(H\to H\)) and \(1\to3\) (\(D\to D\)). With \(F_t\) the freed sets:

\begin{longtable}[]{@{}llll@{}}
\toprule\noalign{}
level & type & free \(\mathcal A\)-slots \(\alpha_j\) & free \(\mathcal B\)-slots \(\beta_j\) \\
\midrule\noalign{}
\endhead
\bottomrule\noalign{}
\endlastfoot
0 & H & \(F_{0\to1} + R_{0\to2}\) & \(R_{4\to0} + F_{3\to0}\) \\
1 & D & \(F_{1\to2} + R_{1\to3}\) & \(F_{0\to1} + F_{4\to1}\) \\
2 & H & \(F_{2\to3} + R_{2\to4}\) & \(F_{1\to2} + R_{0\to2}\) \\
3 & D & \(F_{3\to4} + F_{3\to0}\) & \(F_{2\to3} + R_{1\to3}\) \\
4 & H & \(R_{4\to0} + F_{4\to1}\) & \(F_{3\to4} + R_{2\to4}\) \\
\end{longtable}

(sizes of the sets). With all point sets of size \(f_P\), all line sets of size \(f_L\), and \(m\) elements freed on
each non-friendly transition: \(K_{\mathcal B}\le\min(f_P+m,\ f_L+m,\ 2f_L)\) and
\(K_{\mathcal A}\le\min(f_L+m,\ 2f_P,\ f_P+m)\). Five levels and two matchings form the regular tournament on 5
vertices, and a split of the types into 3 + 2 levels cuts 6 of the 10 arcs, the maximum. The split \(HDHDH\) (the two
\(D\) levels at distance 2) also gives every level at least one friendly outgoing and one friendly incoming transition;
a split with adjacent \(D\) levels, such as \(DDHHH\), does not. (With three levels, some level always has no friendly
outgoing transition, so typing does not help for \(r = 1\).)

\textbf{Seven levels (\(r = 3\)).} The (15,7) graph uses types \(HHDHDHD\): 12 of the 21 transitions are friendly and
every level has at least one friendly transition in each direction (research/d15D7b/README.md {[}b15{]}).

\begin{center}\rule{0.5\linewidth}{0.5pt}\end{center}

\hypertarget{hub-graphs}{%
\subsection{7. Hub graphs}\label{hub-graphs}}

Theorem 3.4 needs a bipartite hub graph \(Y\) of maximum degree \(\le q-r\) (for the degree bound) such that same-colour
hubs are within distance 6 in \(X_R^+\) (condition (HG); paths through the levels may help). This is a purely
combinatorial side condition. Write \(\mathrm{HG}(K_1,K_2;\delta)\) for a bipartite graph with colour classes of sizes
\(K_1, K_2\), maximum degree \(\le\delta\), and all same-colour distances \(\le 6\) \emph{in the graph itself}; if \(Y\)
is such a graph with \(\delta\le q-r\), both conditions hold.

\textbf{Lemma 7.1 (plane with pendant trees).} Let \(s\) be a prime power and \(t, u\ge 0\) integers with
\(s+1+t\le\delta\) and \(u+1\le\delta\). Then \(\mathrm{HG}(K,K;\delta)\) exists for every \(K\) with
\(s^2+s+1\le K\le (s^2+s+1)(1+t+tu)\).

\emph{Proof.} Start from the incidence graph of \(\mathrm{PG}(2,s)\) (the \emph{core}: diameter 3, same-colour distance
2). Attach to some core vertices up to \(t\) pendant vertices each (``depth 1'', of the opposite colour), and to some
depth-1 vertices up to \(u\) pendant vertices each (``depth 2'', of the colour of the core vertex). Degrees are at most
\(s+1+t\), \(u+1\) and \(1\). For two vertices \(x, y\) of the same colour with core ancestors \(c(x), c(y)\),
\(d(x,y)\le \mathrm{depth}(x)+d(c(x),c(y))+\mathrm{depth}(y)\). The core ancestors have the same colour when the depths
have the same parity (core distance \(\le 2\)) and different colours otherwise (core distance \(\le 3\)); the possible
totals are \(0+2+0\), \(0+3+1\), \(0+2+2\), \(1+2+1\), \(1+3+2\), \(2+2+2\), all \(\le 6\). To get \(K\) vertices of
each colour, give \(c\) depth-1 children to the points and \(c\) to the lines, and \(g\) depth-2 children below each of
these two families, with \(c = \lceil (K-s^2-s-1)/(1+u)\rceil\le (s^2+s+1)t\) and \(g = K-s^2-s-1-c\le uc\). \(\square\)

Taking \(t = \delta-s-1\) and \(u = \delta-1\), Lemma 7.1 gives hub graphs with up to
\(\kappa(\delta) := \max_s (s^2+s+1)(1+(\delta-s-1)\delta)\approx \tfrac{4}{27}\delta^4\) vertices per colour (\(s\)
over prime powers \(\le \delta-1\); \(\kappa(3) = 7\), \(\kappa(4) = 35\), \(\kappa(8) = 527\), \(\kappa(9) = 868\),
\(\kappa(10) = 1{,}271\), \(\kappa(11) = 1{,}938\), \(\kappa(14) = 5{,}187\)). For \(\delta\ge 3\) the intervals for
\(s = 2\) and for the optimal \(s\) overlap, so \(\mathrm{HG}(K,K;\delta)\) exists for every
\(7\le K\le\kappa(\delta)\). This is enough for the general statements of Section 9 but far below what the records need.
The record hub graphs were found by computer and checked:

\begin{longtable}[]{@{}
  >{\raggedright\arraybackslash}p{(\columnwidth - 8\tabcolsep) * \real{0.1134}}
  >{\raggedright\arraybackslash}p{(\columnwidth - 8\tabcolsep) * \real{0.1031}}
  >{\raggedright\arraybackslash}p{(\columnwidth - 8\tabcolsep) * \real{0.0619}}
  >{\raggedright\arraybackslash}p{(\columnwidth - 8\tabcolsep) * \real{0.4124}}
  >{\raggedright\arraybackslash}p{(\columnwidth - 8\tabcolsep) * \real{0.3093}}@{}}
\toprule\noalign{}
\begin{minipage}[b]{\linewidth}\raggedright
cell
\end{minipage} & \begin{minipage}[b]{\linewidth}\raggedright
hubs per colour
\end{minipage} & \begin{minipage}[b]{\linewidth}\raggedright
hub degree
\end{minipage} & \begin{minipage}[b]{\linewidth}\raggedright
hub graph
\end{minipage} & \begin{minipage}[b]{\linewidth}\raggedright
(HG) checked by
\end{minipage} \\
\midrule\noalign{}
\endhead
\bottomrule\noalign{}
\endlastfoot
(7,7), (9,7), (10,7), (13,7) & 81, 295, 229, 768 & \(q-1\) & random \((q-1)\)-regular, retried until same-colour
distances \(\le 6\) in \(Y\) & research/d13D7/multi.c (BFS in \(Y\)), then the full all-pairs check \\
(11,7) & 7,300 & 8 & Haar graph of \(\mathbb Z_{73}\rtimes\mathbb Z_{100}\), diameter 7 & build\_haar2.py, then full
check \\
(12,7) & 14,600 & 7 & Haar graph of \((\mathbb Z_{73}\rtimes\mathbb Z_8)\times\mathbb Z_5^2\), with paths through the
levels (``cluster-assisted'', research/squads/cells/d12D7.md §2) & full check \\
(14,7) & 9,014 & 9 & 9 random perfect matchings & BFS in \(Y\), then full check \\
(15,7) & 10,826 / 10,481 & \(\le 8\) & random, repaired by 2-switches; paths through the levels allowed & BFS from all
hubs, then full check \\
(16,7) & 14,884 & \(\le 11\) & voltage graph over \(\mathbb Z_{61}\) & hubcheck.c, then full check \\
\end{longtable}

\begin{center}\rule{0.5\linewidth}{0.5pt}\end{center}

\hypertarget{the-record-graphs-as-corollaries}{%
\subsection{8. The record graphs as corollaries}\label{the-record-graphs-as-corollaries}}

In each case below the graph is \(X_R^+\) of Definition 3.1, every transition is repaired by Theorem B or C, and (HG)
holds, so Theorem 3.4 gives the bound. ``Computer'' marks the parts that are found by search and checked by computer
rather than proved in general: the specific sets, the bijection \(f\), and the hub graph. The theorems themselves are
proved for all \(q\).

\textbf{Corollary 8.1 (three levels, one local duality: (7,7), (9,7), (10,7), (13,7)).} Let \(r = 1\), all types \(H\),
\(M_t = \sigma\) for all three transitions and \(R_t = R\) for all three, with \((R,\sigma)\) as in Proposition 5.2 /
Corollary 5.3(b). If the hub graph \(Y\) (with \(\lvert R\rvert\) hubs of each colour) has maximum degree \(\le q-1\)
and satisfies (HG) of Theorem 3.4 (for instance, if it is an \(\mathrm{HG}(\lvert R\rvert,\lvert R\rvert;q-1)\)), then
\[n_b(q+2,\,7)\ \ge\ 6N(q) + 2\lvert R\rvert .\] Records: \(q = 5\): \(\lvert R\rvert = 81\) (\(R_0\) = 72 pairwise
opposite points, 9 extra points, spread lines of \(q^3+1 = 126\) available), \(n = 23{,}598\); \(q = 7\): \(288 + 7\),
\(n = 118{,}238\); \(q = 8\): \(217 + 12\), \(n = 225{,}152\); \(q = 11\): \(756 + 12\), \(n = 1{,}064{,}472\).
Computer: \(R_0\) (for \(q = 7, 11\) an orbit of the Weyl group \(W(E_7)\) acting on \(Q(6,q)\) through the \(E_7\) root
lattice mod \(q\), research/cell147/ovoid/ and octo/ {[}b14{]}; being pairwise non-orthogonal in \(Q(6,q)\) it is a
partial ovoid of \(H(q)\) by (F4); for \(q = 5, 8\) from research/d13D7/q57 and q8 {[}b14{]}), the extra points \(X\),
the lines \(L_x\), the bijection \(f\) onto Hermitian spread lines, the hub graph.

\textbf{Corollary 8.2 (duality, \(q = 3^h\): (11,7), (12,7)).} Let \(q = 3^h\), \(\varphi\) a duality, all types \(H\),
\(M_t = \varphi\vert_P\) for every transition, and \(R_t = F\) for every transition, \(F\) the point set of a spread
(e.g.~the Hermitian spread), \(\lvert F\rvert = (q+1)(q^3+1)\). By Theorem C every transition is repaired. Every level
has \(r\lvert F\rvert\) free slots of each colour. If \(K\le r(q+1)(q^3+1)\) hubs of each colour are attached and the
hub graph \(Y\) has maximum degree \(\le q-r\) and satisfies (HG) (for instance, if it is an \(\mathrm{HG}(K,K;q-r)\)),
then \[n_b(q+1+r,\,7)\ \ge\ (2r+1)\cdot 2N(q) + 2K .\] Records (\(q = 9\), \(\lvert F\rvert = 7{,}300\)): \(r = 1\),
\(K = 7{,}300\): \(n = 413{,}180\); \(r = 2\), \(K = 14{,}600\): \(n = 693{,}500\). By Proposition 1.4 and Remark 4.2 no
larger freed set leaves every transition repaired, so within this construction, and as long as condition (T) of Theorem
3.4 is required, \(K = r(q+1)(q^3+1)\) is the cap. Computer: the duality \(\varphi\) of \(H(9)\) (explicit, verified),
the hub graphs. Proved: the freed set and its size, for every \(q = 3^h\).

\textbf{Corollary 8.3 (typed levels: (14,7), (15,7), (16,7)).} Let the types be \(HDHDH\) (\(r = 2\)) or \(HHDHDHD\)
(\(r = 3\)), let every friendly transition carry \(M_t = \mathrm{id}\) and a segment set (points) or dual segment set
(lines) \(F_t\), and let the non-friendly transitions carry \(M_t = \sigma\) (\(H\to H\)) or \(\sigma^{-1}\)
(\(D\to D\)) with freed sets \(R\) (resp. \(\sigma(R)\)) satisfying Theorem B. If \(K_{\mathcal B}\le\min\alpha_j\),
\(K_{\mathcal A}\le\min\beta_j\) (Section 6) and the hub graph \(Y\) has maximum degree \(\le q-r\) and satisfies (HG)
(for instance, if it is an \(\mathrm{HG}(K_{\mathcal B},K_{\mathcal A};q-r)\)), then
\[n_b(q+1+r,\,7)\ \ge\ (2r+1)\cdot 2N(q) + K_{\mathcal A} + K_{\mathcal B}.\] Records:

\begin{itemize}
\tightlist
\item
  (14,7), \(q = 11\), \(HDHDH\): all friendly point sets equal one non-collinear set of 8,645 points, all friendly line
  sets one non-concurrent set of 8,246 lines (both unions of orbits of an order-133 collineation), and the 768-point
  mirrored set of Corollary 8.1 on the four non-friendly transitions. \(\alpha = (9413, 9014, 9413, 16492, 9413)\),
  \(\beta = (9014, 17290, 9014, 9413, 9014)\), \(K_{\mathcal A} = K_{\mathcal B} = 9{,}014\): \(n = 1{,}789{,}588\).
\item
  (15,7), \(q = 11\), \(HHDHDHD\): twelve friendly sets with 5,674--9,303 elements each; in the four bottleneck sets
  some elements have one partner (5,577 pairs in all), so these are segment sets but not non-collinear; the 768-point
  mirrored set on the nine non-friendly transitions. \(K_{\mathcal B} = 10{,}826\), \(K_{\mathcal A} = 10{,}481\):
  \(n = 2{,}501{,}491\). (The graph also satisfies the lab's unnecessary hub condition (R2).)
\item
  (16,7), \(q = 13\), \(HDHDH\): a non-collinear set of 14,762 points and a non-concurrent set of 14,823 lines (unions
  of orbits of an order-61 element of \(SL(3,13)\)), and on the four non-friendly transitions a pairwise opposite set of
  122 points (resp. lines) with \(\sigma\) as in Corollary 5.3(a).
  \(K_{\mathcal A} = K_{\mathcal B} = 14{,}762 + 122 = 14{,}884\): \(n = 4{,}052{,}108\).
\end{itemize}

Computer: all these sets, \(\sigma\), and the hub graphs.

\textbf{8.4 Where the exact graphs are.} Each record graph is determined by a solution file (research/frontier/sols/ on
master) that lists the freed sets, the matchings (or the data that determine them), the slot \(\omega_j(h)\) of every
hub in every level, and the hub graph \(Y\); the builder named here turns it into an explicit edge list (which the
coordinator then checked, Section 10.2). So the neighbourhood of every vertex, including every hub, can be read off.

\begin{longtable}[]{@{}lll@{}}
\toprule\noalign{}
cell & solution file (research/frontier/sols/) & builder \\
\midrule\noalign{}
\endhead
\bottomrule\noalign{}
\endlastfoot
(7,7) & d7D7\_hex5\_mirror81\_n23598.sol & research/d13D7/tools/multi\_sol.py \\
(9,7) & d9D7\_hex7\_mirror295\_n118238.sol & same \\
(10,7) & d10D7\_hex8\_mirror229\_n225152.sol & same \\
(11,7) & d11D7\_hex9\_spread730\_haarMeta73x100\_n413180.sol, .hubs & research/d10D7/dual/xgen.c \\
(12,7) & d12D7\_hex9\_spread730\_clhaarZ73Z8xZ5sq\_n693500.sol, .hubs3 & research/squads/d12D7/clhaar/x5gen3.c \\
(13,7) & d13D7\_hex11\_mirror768\_n1064472.sol & research/d13D7/tools/multi\_sol.py \\
(14,7) & d14D7\_hex11\_typedmirror768\_n1789588.sol & research/cell147/typedmirror/tm\_build.py \\
(15,7) & d15D7\_hex11\_typedrelax2\_n2501491.sol & research/d15D7b/tm3\_build.py \\
(16,7) & d16D7\_hex13\_typed61\_K244\_n4052108.sol & research/squads/cells/d16D7/typed5s\_build.py \\
\end{longtable}

\begin{center}\rule{0.5\linewidth}{0.5pt}\end{center}

\hypertarget{what-is-proved-for-every-q}{%
\subsection{\texorpdfstring{9. What is proved for every
\(q\)}{9. What is proved for every q}}\label{what-is-proved-for-every-q}}

Theorems A, B, C and 3.4 hold for every prime power \(q\) (the proofs only use (F1)--(F3), and at least two points on a
line and two lines on a point). What is not known in general is how large the sets can be made and whether large hub
graphs exist; Lemma 1.2 and Lemma 7.1 give explicit choices. With \(\kappa(\delta)\) as in Section 7:

\textbf{Corollary 9.1 (every \(q\ge 4\), three levels).} With \(R\) the \(q^2+1\) points of an elliptic quadric
\(Q^-(3,q)\) section (Lemma 1.2(a)) and \(R^*\) any \(q^2+1\) lines of the Hermitian spread (Corollary 5.3(a)):
\[n_b(q+2,\,7)\ \ge\ 6N(q) + 2\min\bigl(q^2+1,\ \kappa(q-1)\bigr),\] which is \(6N(q)+2(q^2+1)\) for every \(q\ge 5\).

\textbf{Corollary 9.2 (\(q = 3^h\)).} For every \(r\) with \(1\le r\le q-3\) (Corollary 8.2 with Lemma 7.1; here
\(\kappa(q-r) < q^4 < r(q+1)(q^3+1)\), so the hub graph is the only limit):
\[n_b(q+1+r,\,7)\ \ge\ (2r+1)\cdot 2N(q) + 2\kappa(q-r).\]

\textbf{Corollary 9.3 (every \(q\ge 5\), typed levels \(HDHDH\)).} Friendly point sets: the point set of the Hermitian
spread (\((q+1)(q^3+1)\) points); friendly line sets: the lines through the points of an elliptic quadric section
(\((q+1)(q^2+1)\) lines, a dual segment set); non-friendly: \(m = q^2+1\) as in Corollary 9.1. By the table of Section 6
the smallest slot counts are \(\min_j\alpha_j = \min_j\beta_j = (q+1)(q^2+1) + q^2+1 = (q+2)(q^2+1)\), so
\[n_b(q+3,\,7)\ \ge\ 10N(q) + 2\min\bigl((q+2)(q^2+1),\ \kappa(q-2)\bigr).\] The minimum is \(\kappa(q-2)\) for
\(q\le 13\) and \((q+2)(q^2+1)\) for \(q\ge 16\).

The record values exceed these because the sets found by computer are of order \(q^4\) (non-collinear sets of 8,645 and
14,762 points at \(q = 11, 13\), against \(q^4 = 14{,}641\), \(28{,}561\)) and the hub graphs found by computer are much
larger than \(\kappa\). For \(q = 3^h\) the sets are already optimal for Theorem C (Proposition 1.4) and only the hub
graphs are found by computer.

\textbf{Remark 9.4 (consequences of Theorem C not yet built).} Theorem C allows friendly sets that are much larger than
the non-collinear and non-concurrent sets used in the (14,7)--(16,7) records: up to \((q+1)(q^3+1)\) (Proposition 1.4),
e.g. the 15,984 points of the Hermitian spread of \(H(11)\) against 8,645, and the \(12\cdot 756 = 9{,}072\) lines
through the 756-point partial ovoid of \(H(11)\) against 8,246 (non-concurrent). In the \(HDHDH\) layout with these sets
and the 768 mirrored elements the slot counts of Section 6 give
\(K_{\mathcal A} = K_{\mathcal B} = 9{,}072+768 = 9{,}840\), i.e. \(n = 1{,}771{,}560 + 19{,}680 = 1{,}791{,}240\) at
(14,7), provided a hub graph with 9,840 hubs per colour and degree 9 exists. We have not built or checked these graphs;
they are left to the lab's search sessions.

\begin{center}\rule{0.5\linewidth}{0.5pt}\end{center}

\hypertarget{computer-checks}{%
\subsection{10. Computer checks}\label{computer-checks}}

\hypertarget{small-q-tests-of-every-theorem}{%
\subsubsection{\texorpdfstring{10.1 Small-\(q\) tests of every
theorem}{10.1 Small-q tests of every theorem}}\label{small-q-tests-of-every-theorem}}

Two independent sets of programs (details, commands and logs in the two README files):

\begin{itemize}
\tightlist
\item
  research/theorems/d7hex\_tests/: builds \(H(q)\) from scratch from the coordinates of Section 1.1 (\(q = 2,3,4,5\)),
  builds each graph from Definitions 2.1 and 3.1 (asserting that every matching is a bijection), and checks it with
  research/lifts/coordinator\_check/apsp\_mt.c (breadth-first search from every vertex; ``VALID'' = bipartite, maximum
  degree \(\le\Delta\), diameter \(\le 7\));
\item
  research/theorems/d7hex\_lemma41/: uses the lab's generator research/d10D7/hexgen2.c and duality finder
  research/d10D7/dual/dualfind.c; checks Lemma 4.1 directly on all ordered pairs \((u,v)\), and the duality
  construction.
\end{itemize}

\begin{longtable}[]{@{}
  >{\raggedright\arraybackslash}p{(\columnwidth - 6\tabcolsep) * \real{0.1373}}
  >{\raggedright\arraybackslash}p{(\columnwidth - 6\tabcolsep) * \real{0.5098}}
  >{\raggedright\arraybackslash}p{(\columnwidth - 6\tabcolsep) * \real{0.0588}}
  >{\raggedright\arraybackslash}p{(\columnwidth - 6\tabcolsep) * \real{0.2941}}@{}}
\toprule\noalign{}
\begin{minipage}[b]{\linewidth}\raggedright
statement
\end{minipage} & \begin{minipage}[b]{\linewidth}\raggedright
what was built
\end{minipage} & \begin{minipage}[b]{\linewidth}\raggedright
\(q\)
\end{minipage} & \begin{minipage}[b]{\linewidth}\raggedright
result
\end{minipage} \\
\midrule\noalign{}
\endhead
\bottomrule\noalign{}
\endlastfoot
§1.1 coordinates, (F1), (F4), Lemma 1.2(b) & \(H(q)\) from Tits' equations; counts, diameter 6, girth 12 (distance
profile), (F4) on all pairs (\(q\le 4\)) or a sample (\(q = 5\)); Hermitian spread in \(x_0 = c x_4 + e x_3\) & 2, 3, 4,
5 & all hold; spreads of 28, 65, 126 pairwise opposite lines \\
Theorem A & plain level graphs with random bijections, types \(HHH\), \(HDH\), \(HHHHH\), \(HDHDH\) (and \(HHDHHDD\),
\(HHHHHHH\) for \(q = 3, 4\)) & 3, 4, 5 & all VALID, \((q+1+r)\)-regular \\
Lemma 4.1 & all ordered pairs, for spread-point sets, three random maximal segment sets, a random maximal (R1) set;
points and lines & 3, 4, 5 & 0 failures; each set plus one turning point: 36--320 failing pairs \\
Theorem C (friendly) & two levels \(H\to D\), \(D\to H\) and five levels \(HDHDH\), main graph without hubs, with: all
points of the Hermitian spread (112, 325, 756 \(=(q+1)(q^3+1)\)), random maximal segment sets, (R1) sets with no hub
condition, pencils of a partial ovoid, dual segment sets; 3 seeds & 3, 4, 5 & 108 of 108 VALID; set plus one turning
point: 72--1,800 ordered pairs at distance \(>7\) (apsp\_mt: diameter 8, INVALID) \\
Theorem C (friendly, larger \(q\)) & two levels \(H\to D\), all \((q+1)(q^3+1)\) points of the Hermitian spread freed
(2,752 and 15,984 points), main graph & 7, 11 & VALID (\(n\) = 78,432 and 708,624; all sources) \\
Theorem C (duality), Cor. 8.2 & all levels \(H\), every transition matched by a polarity \(\varphi\) of \(H(3)\), the
112 spread points freed on every transition, \(r = 1, 2\) & 3 & VALID (\(n = 2{,}184\) and \(3{,}640\)); one more point
freed: diameter 8 \\
Theorem B, Prop. 5.2, Cor. 5.3 & \(M\) built exactly as in Proposition 5.2, (H1)--(H3) re-checked on it; \(R\) pairwise
opposite (an ovoid of 28 points at \(q = 3\); 23; 33) and mirrored \(R\) with pairs at distance 4 (17 points/19 pairs;
26/21; 38/28); types \(HHH\) (same \(M\), \(R\) on all transitions) and \(HDHDH\) (\(M\) on \(H\to H\), \(M^{-1}\) on
\(D\to D\)) & 3, 4, 5 & all VALID; breaking (H3) at one point: 50, 76, 126 ordered pairs at distance \(>7\) \\
Theorem 3.4 (hubs) & full graphs with hubs on random free slots and a random hub graph of degree \(\le q-r\): \(HHH\),
\(HDH\) (\(K\) = slot minimum, 28 and 35), \(HDHDH\) (\(K\) = 48 and 61, below the slot minimum because the hub degree
\(q-2\) is small) & 4, 5 & all VALID (\(n\) = 8,246, 8,246, 13,746; 23,506, 23,506, 39,182) \\
\end{longtable}

Every positive test passed and every negative control failed. Small \(q\) cannot test the hub graphs at the record sizes
(at \(q\le 5\) the hub degree \(q-r\) is 1 to 4); for those see Section 10.2.

\hypertarget{the-record-graphs}{%
\subsubsection{10.2 The record graphs}\label{the-record-graphs}}

Each record graph was rebuilt by the lab's coordinator (on a separate machine) from its solution file with the session's
builder, which also re-checks the hypotheses ((H1)--(H3) at every freed element, non-collinearity or the partner
condition of the friendly sets, degrees, the hub graph). The edge list was then checked by
research/lifts/coordinator\_check/apsp\_mt.c: simple graph, maximum degree \(\le\Delta\), bipartite, connected, and a
bit-parallel breadth-first search from \textbf{every} vertex, no symmetry assumed. All nine runs are in
research/coordinator\_check/check\_2026-10-25.txt:

\begin{longtable}[]{@{}lrl@{}}
\toprule\noalign{}
cell & \(n\) & apsp\_mt output (all sources) \\
\midrule\noalign{}
\endhead
\bottomrule\noalign{}
\endlastfoot
(7,7) & 23,598 & RESULT VALID d=7 D=7 n=23598 \\
(9,7) & 118,238 & RESULT VALID d=9 D=7 n=118238 \\
(10,7) & 225,152 & RESULT VALID d=10 D=7 n=225152 \\
(11,7) & 413,180 & RESULT VALID d=11 D=7 n=413180 \\
(12,7) & 693,500 & RESULT VALID d=12 D=7 n=693500 \\
(13,7) & 1,064,472 & RESULT VALID d=13 D=7 n=1064472 \\
(14,7) & 1,789,588 & RESULT VALID d=14 D=7 n=1789588 \\
(15,7) & 2,501,491 & RESULT VALID d=15 D=7 n=2501491 \\
(16,7) & 4,052,108 & RESULT VALID d=16 D=7 n=4052108 \\
\end{longtable}

(The finders' own checks were partly reduced by symmetry, e.g.~an order-61 automorphism for (16,7); the coordinator's
runs above assume none.)

\begin{center}\rule{0.5\linewidth}{0.5pt}\end{center}

\hypertarget{prior-art-and-what-is-new}{%
\subsection{11. Prior art, and what is new}\label{prior-art-and-what-is-new}}

\emph{Classical.} Generalized hexagons and their incidence graphs {[}Ti59{]}; the split Cayley hexagon in \(Q(6,q)\),
its dualities for \(q = 3^h\), and the Hermitian spread ({[}Ti59{]}, {[}Th80{]}, {[}VM98{]}).

The level graph of Definition 2.1 is a bipartite analogue of the \emph{star product} of Bermond, Delorme and Farhi
{[}BDF82, BDF84{]} (see {[}MS13, §3.3.2{]}): copies of \(I\) indexed by the vertices of the complete graph \(K_{2r+1}\),
oriented as a regular tournament, and for each arc a perfect matching between one colour class of the tail copy and the
other colour class of the head copy. The survey credits the constructions of large bipartite graphs mainly to Delorme's
bipartite versions of the star product and of compounding {[}De85a, De85b{]}, to the partial Cayley graphs of Bond and
Delorme, and to the voltage lifts of Loz, Pérez-Rosés and Pineda-Villavicencio, which updated about a third of the table
{[}MS13, §3.4.5{]}. The D = 7 values with \(\Delta\ge 9\), and (7,7), have exactly the form \((2r+1)\cdot 2N(q)\) of
Theorem A; the survey's table does not name the source of each cell, and we have not determined it. We do not claim
Theorem A as new: it is the routine argument for such level graphs.

\emph{New here (as far as we know).} Deleting matching edges at freed elements and reattaching the freed slots to hub
vertices (Definition 3.1, Theorem 3.4); the repair theorems B and C; the use of dual levels (``typed levels'') to make
transitions automorphic; segment sets and Proposition 1.4; mirrored sets (Propositions 5.1, 5.2); and the nine values in
Section 0.

\hypertarget{appendix-notes-for-readers-of-the-lab-notes}{%
\subsection{12. Appendix: notes for readers of the lab notes}\label{appendix-notes-for-readers-of-the-lab-notes}}

\emph{What is new relative to the per-cell notes.} Writing the proofs out gave one simplification: Theorem C replaces
the lab's conditions on freed sets at friendly transitions (non-collinear sets, and the ``relaxed sets'' with hypotheses
(R1), (R2) of research/d15D7b/RELAXED.md) by the single, weaker condition of being a segment set. The lab's (12,7) notes
(research/squads/cells/d12D7.md §2) state the same condition for transitions matched by a duality and check it by
computer at \(q = 9\); Theorem C proves it and transfers it to friendly transitions. The record graphs are unaffected
(they satisfy the stronger hypotheses and were checked by computer in full), but larger sets are now allowed (Remark
9.4).

\emph{Why the lab's (R2) was not needed.} The lab proof of the relaxed theorem (RELAXED.md, cases 4 and 5) only used
paths whose matching edge sits at a point collinear with the freed point \(r\). When \(r\) has a partner \(x\) and the
target is a line \(u\) with \(d(r,u) = 5\) and geodesic \(r - y - x - y_2 - w_2 - u\), the only such point is \(x\)
(freed), and (R2) supplies a hub. The path through \(w_2\) (at distance 4 from \(r\)) avoids this: \(w_2\) is collinear
with \(x\), so it is not freed (else \(x\) would have two partners), and \((r,j')\leadsto(w_2,j')\to(w_2,j)\to(u,j)\)
has length \(4+1+1 = 6\). This is the first step of the proof of Lemma 4.1 (a point and a line at distance 5: three
points in a row on the geodesic).

\emph{Hubs in Theorem B.} The earlier version of Theorem B (research/d13D7/REPORT.md §2, with \(R\) pairwise opposite)
routed case 6 through the hub of \(x\) in two levels. The route of case 6 in the proof of Theorem B uses no hub, which
is what lets hubs take their slots from different freed elements in different levels (Definition 3.1).

\hypertarget{references}{%
\subsection{References}\label{references}}

\begin{itemize}
\tightlist
\item
  {[}BDF82{]} J.-C. Bermond, C. Delorme, G. Farhi, Large graphs with given degree and diameter III, Ann. Discrete Math.
  13 (1982) 23--32.
\item
  {[}BDF84{]} J.-C. Bermond, C. Delorme, G. Farhi, Large graphs with given degree and diameter II, J. Combin. Theory
  Ser. B 36 (1984) 32--48.
\item
  {[}BGS06{]} S. Ball, P. Govaerts, L. Storme, On ovoids of parabolic quadrics, Des. Codes Cryptogr. 38 (2006) 131--145.
\item
  {[}De85a{]} C. Delorme, Grands graphes de degré et diamètre donnés, European J. Combin. 6 (1985) 291--302.
\item
  {[}De85b{]} C. Delorme, Large bipartite graphs with given degree and diameter, J. Graph Theory 9 (1985) 325--334.
\item
  {[}MS13{]} M. Miller, J. Širáň, Moore graphs and beyond: a survey of the degree/diameter problem, Electron. J.
  Combin., Dynamic Survey DS14 (2nd ed., 2013).
\item
  {[}Th80{]} J. A. Thas, Polar spaces, generalized hexagons and perfect codes, J. Combin. Theory Ser. A 29 (1980)
  87--93.
\item
  {[}Th81{]} J. A. Thas, Ovoids and spreads of finite classical polar spaces, Geom. Dedicata 10 (1981) 135--143.
\item
  {[}Ti59{]} J. Tits, Sur la trialité et certains groupes qui s'en déduisent, Publ. Math. IHÉS 2 (1959) 13--60.
\item
  {[}VM98{]} H. Van Maldeghem, Generalized Polygons, Birkhäuser, 1998.
\end{itemize}

\end{document}
